\chapter{Réécritures pour la stabilité numérique}
\label{ch:stabilite}

\begin{lecture}
Les formules démontrées jusqu'ici sont exactes en arithmétique réelle. Sur une
machine, les nombres ont une taille finie, et deux expressions mathématiquement
égales n'y sont pas équivalentes : l'une rend un résultat, l'autre échoue. Les
deux réécritures établies ici conservent chaque valeur et modifient seulement la
suite d'opérations qui y conduit.
\end{lecture}

\section{Coût exprimé à partir du score}

En double précision, $\sigma(z)$ vaut exactement $1$ dès que $z$ dépasse environ
$37$ : $e^{-37}$ est de l'ordre de $10^{-17}$, en dessous de ce qu'un flottant
distingue de $1$ lorsqu'on l'y ajoute. Le coût demande alors $\log(1-1)$, qui
n'existe pas. Un élève d'une autre maison auquel le modèle attribue une forte
probabilité d'être \Gryf suffit donc à interrompre le programme, alors que
c'est exactement le cas que le coût doit signaler.

\begin{rappel}[Logarithme d'un quotient, logarithme d'une exponentielle]
$\log(a/b)=\log a-\log b$, et $\log(e^{b})=b$.
\end{rappel}

Reprenons le coût d'un élève et remplaçons $p$ par $\sigma(z)$ :

\begin{align}
  \ell &= -\bigl[y\log\sigma(z)+(1-y)\log(1-\sigma(z))\bigr] \notag\\[4pt]
  \log\sigma(z) &= \log\frac{1}{1+e^{-z}}
    = \log 1-\log(1+e^{-z}) \notag\\
    &= -\log(1+e^{-z})
    &&\text{car } \log 1=0 \notag\\[4pt]
  \log\bigl(1-\sigma(z)\bigr) &= \log\frac{e^{-z}}{1+e^{-z}}
    = \log e^{-z}-\log(1+e^{-z}) \notag\\
    &= -z-\log(1+e^{-z})
    &&\text{car } \log e^{-z}=-z \notag\\[4pt]
  \ell &= -\Bigl[y\bigl(-\log(1+e^{-z})\bigr) \notag\\
    &\qquad\quad +(1-y)\bigl(-z-\log(1+e^{-z})\bigr)\Bigr] \notag\\[4pt]
  &= y\log(1+e^{-z})+(1-y)z+(1-y)\log(1+e^{-z})
    &&\text{on distribue} \notag\\[4pt]
  &= \log(1+e^{-z})+(1-y)z
    &&\text{les logs se regroupent} \notag\\[4pt]
  &= \underbrace{\log(1+e^{-z})+z}_{=\;\log(1+e^{z})}-yz \notag\\[4pt]
  &= \log(1+e^{z})-yz
\end{align}

Le regroupement souligné s'établit ainsi :
\[
  \log(1+e^{-z})+z
  = \log(1+e^{-z})+\log e^{z}
  = \log\bigl(e^{z}(1+e^{-z})\bigr)
  = \log(e^{z}+1),
\]
en remontant $z$ sous forme de $\log e^{z}$, en regroupant les deux logarithmes
en celui d'un produit, puis en développant : $e^{z}\cdot e^{-z}=e^{0}=1$.

\begin{resultat}[Forme employée par le programme]
\[
  \ell(y,z)=\operatorname{softplus}(z)-yz,
  \qquad \operatorname{softplus}(z)=\log(1+e^{z}).
\]
Cette forme se calcule à partir du score seul, hors d'atteinte des arrondis qui
produisent $\log 0$.
\end{resultat}

\begin{releve}
Pour $y=1$ et $z=6$, la forme directe donne $-\log(0{,}997527)=0{,}002476$ et
celle-ci $\operatorname{softplus}(6)-6=6{,}002476-6=0{,}002476$.
\end{releve}

\section{Formes protégées de softplus et de la sigmoïde}

$\operatorname{softplus}$ contient encore $e^{z}$, qui déborde pour $z$ grand.
On l'écrit donc
\[
  \operatorname{softplus}(z)=\max(z,0)+\log\bigl(1+e^{-|z|}\bigr),
\]
dont l'exposant est toujours négatif ou nul. La vérification porte sur deux cas.
Pour $z\geq0$, $\max(z,0)=z$ et l'expression vaut
$z+\log(1+e^{-z})=\log(e^{z}+1)$, par le même regroupement que ci-dessus. Pour
$z<0$, $\max(z,0)=0$ et $|z|=-z$, il reste $\log(1+e^{z})$, qui est la
définition.

La sigmoïde reçoit le même traitement, avec ses deux écritures
(ch.~\ref{ch:logit}) :
\[
  \sigma(z)=
  \begin{cases}
    1/(1+e^{-z}), & z\geq 0,\\[2pt]
    e^{z}/(1+e^{z}), & z<0 .
  \end{cases}
\]
Dans les deux branches, l'exponentielle porte sur un nombre négatif ou nul. Un
exposant très négatif s'arrondit à zéro sans dommage ; un exposant très positif dépasse en revanche la capacité du format et provoque une erreur.

\begin{figure}[htbp]
\centering
\begin{tikzpicture}[font=\small,x=0.0083cm]
% Axe des scores, de -900 a 900 ; 1 unite = 0,0083 cm, soit 14,9 cm au total.
\fill[warning!18] (-900,-0.28) rectangle (-710,0.28);
\fill[warning!18] (710,-0.28) rectangle (900,0.28);
\fill[huffle!12] (37,-0.28) rectangle (710,0.28);
\fill[huffle!12] (-710,-0.28) rectangle (-37,0.28);
\fill[rulecolor!12] (-37,-0.28) rectangle (37,0.28);
\draw[ink!45] (-900,0) -- (900,0);
\foreach \v in {-800,-400,0,400,800}{
  \draw[ink!45] (\v,-0.32) -- (\v,0.32);
  \node[font=\scriptsize,text=ink,anchor=north] at (\v,-0.36) {$\v$};
}
\foreach \v/\lab in {-710/{$-710$},-37/{$-37$},37/{$37$},710/{$710$}}{
  \draw[ink!60,thin] (\v,-0.28) -- (\v,0.28);
  \node[font=\scriptsize,text=ink!75,anchor=south] at (\v,0.32) {\lab};
}
\node[font=\footnotesize,text=ink,anchor=north] at (0,-1.0) {score $z$};
\node[font=\scriptsize,text=rulecolor,align=center,anchor=north]
  at (0,-1.6) {$\sigma$ et $\log$\\ exacts};
\node[font=\scriptsize,text=huffle,align=center,anchor=north]
  at (370,-1.6) {$\sigma$ arrondie à $0$ ou $1$ :\\ $\log$ direct impossible};
\node[font=\scriptsize,text=huffle,align=center,anchor=north]
  at (-370,-1.6) {$\sigma$ arrondie à $0$ ou $1$ :\\ $\log$ direct impossible};
\node[font=\scriptsize,text=warning,align=center,anchor=north]
  at (805,-1.6) {$e^{|z|}$\\ déborde};
\node[font=\scriptsize,text=warning,align=center,anchor=north]
  at (-805,-1.6) {$e^{|z|}$\\ déborde};
\end{tikzpicture}
\caption{Les trois régimes du calcul en double précision. La zone centrale seule
tolère l'écriture directe ; les formes protégées rendent des valeurs finies sur
tout l'axe. Un score de $\pm800$ apparaît dès que les coefficients divergent
(ch.~\ref{ch:convexite}).}
\label{fig:regimes}
\end{figure}
\FloatBarrier

\begin{table}[htbp]
\centering
\begin{tabular}{rll}
\toprule
$z$ & écriture directe & écriture protégée \\
\midrule
$-800$ & $e^{800}$ déborde & $0{,}0$ \\
$+800$ & $\log(1+e^{800})$ déborde & $800{,}0$ \\
$+40$  & $\log(1-\sigma(40))=\log 0$ & $0{,}0$ \\
\bottomrule
\end{tabular}
\caption{En double précision, l'exponentielle déborde au-delà de
$|z|\approx710$ et la sigmoïde s'arrondit à $1$ dès $z\approx37$. Les formes
protégées rendent des valeurs finies dans les trois cas.}
\label{tab:stabilite}
\end{table}

\begin{vigilance}[Contrôles numériques]
Le premier coût vaut $\log 2=0{,}693147$ quand les coefficients partent de zéro.
Toute autre valeur signale une erreur dans le score ou l'initialisation.

Le gradient calculé par la formule, comparé à
$\bigl(J(\w+h\mathbf e_j)-J(\w-h\mathbf e_j)\bigr)/2h$ sur un cas de petite
taille, doit coïncider à $10^{-6}$ près. Ce test détecte un signe inversé, un
facteur $1/m$ oublié, un indice décalé.

La sigmoïde et le coût, évalués en $z=\pm1000$, doivent rendre des nombres finis.
\end{vigilance}
